\section{Search superdiffusion}\label{sec:search_superdiffusion}
Superdiffusion is often a prominent pattern in natural dynamical systems, and its presence points to a self-similar fractal structure. Such dynamics have not been the focus of other studies, and visual trajectories have often been treated as fixational Brownian motion, externally interrupted by saccades, which would suggest separate regimes of diffusion across $q$. We present here a monoscaling superdiffusive pattern in gaze dynamics across a wide range of scales. Many studies that have described eye movement dynamics as normally diffusive stem from the sensitivity of data-processing choices in detecting anomalous diffusion at confined boundaries, such as fixations. Conventional analyses often inadvertently obscure critical spatiotemporal features.

The most important methodological choice in accurately characterising visual diffusion patterns is to time-average displacements. As demonstrated by \textcite{vilk_strong_2025}, aggregating data across disparate temporal regimes can wash out localised scaling behaviours, yielding normal diffusion as an artefact. Studies that construct regimes of normality and subdiffusive regimes are thus common in the study of eye movements, but this ignores continuous-time superdiffusion, which we unify across multiple scales. For such studies, see \cite{lencastre_dynamical_2025, leroux_random_2021, roberts_fixational_2013, collewijn_significance_2008}.

\subsubsection*{Analysis of multiscale superdiffusion in eye-tracking data}
We analyse the binocular-averaged gaze trajectories over all 22 human subjects provided from the study in \textcite{naher_primate_saccades}, as described in \cref{sec:model_setup}. From the onset of the search task, we time-average the fractional absolute displacements $\langle \|\Delta \textbf{x}(\Delta t) \|^q \rangle $ for 30 values of $q$ linearly separated from $q=0.1$ to $q=3.0$. These fractional absolute moments are evaluated over all valid time lags $\Delta t < \qty{1000}{\ms}$ with a resolution of $\Delta t = \qty{1}{\ms}$. All intervals containing NaN values (such as blinks) are discarded. We then ensemble-average these displacements across all subjects. These results are shown in \cref{fig:search_superdiffusion}\textbf{a-d)}.

We perform a linear best-fit of the gaze and model fractional moments in log-log space $M_q(\Delta t)$ as defined in \cref{eq:Mq} to find short-time power-law diffusive scaling with time, thus sampling from our moment spectrum function $\lambda(q)$ as defined in \cref{eq:MFD}. Each of these fractional moments is overlaid in \cref{fig:search_superdiffusion}\textbf{a,c)}, and their slopes are plotted in \cref{fig:search_superdiffusion}\textbf{b,d)}. The moment spectrum function $\lambda(q)$ is thus found to have a good fit to $\lambda(q) = (0.897 \pm 0.007)q$ with an error in the fit of $R^2 = 0.999$. This behaviour yields monoscaling across spatial scales, from fixations to saccades. The corresponding Lévy exponent given such linear monoscaling in $\lambda(q)$ is thus found to be $\alpha = 1.115 \pm 0.007$.
\subsection{Spectrum of fractional moments}\label{subsec:spectrum_of_fractional_moments}
\Cref{fig:search_superdiffusion} summarises the scaling of the time-averaged fractional absolute displacements. For comparison of the various scaling exponents, we define a rescaled moment function $M_q(\Delta t)$ to give a good comparison of the spectrum of moments under the same dimensions of distance. $M_q(\Delta t)$ is here defined
\begin{equation} \label{eq:Mq}
    M_q(\Delta t) = C_q \langle \|\Delta\mathbf{x}(\Delta t)\|^{q} \rangle^{1/q} \sim \Delta t^{\lambda(q) / q}
\end{equation}
with $C_q$ defined such that $M_q(\Delta t_{\text{min}})=1$ to give a shared origin at the first time lag $\Delta t$. Small $q$ gives strong weight to the small, frequent displacements of fixational motion, while large $q$ weights the rare, large displacements which are given in saccades. The shape of $\lambda(q)$, therefore, examines the spectrum of exponents that may describe superdiffusion at any scale. A linear spectrum, $\lambda(q) = q/\alpha$, is monoscaling: all moments grow with the same exponent, as for a self-similar Lévy process with stability index $\alpha$. A nonlinear spectrum is the general case found in many systems called strong anomalous diffusion, which consists of two dynamical strategies stitched together at different scales \cite{vilk_strong_2025, cagnetta_strong_2015, castiglione_strong_1999, wang_strong_2020}. See \cref{subsec:levy_motion} for a further discussion of mono- and multiscaling, and \cref{fig:SANDvsmono} for its schematic representation.

\Cref{fig:search_superdiffusion} shows the rescaled moments $M_q(\Delta t)$ for the human eye-tracking data and modelled gaze. The slope of each curve in log--log space is robustly linear, with constant slope $\lambda(q)/q$. Lévy monoscaling predicts a slope across the entire spectrum, with $\lambda(q) = q/\alpha$. For comparison, Brownian motion ($\alpha = 2$) has slope $1/2$ and ballistic motion has slope $1$ across the entire spectrum $q$. The monoscaling exponent is found by performing a linear regression on our spectrum $\lambda(q)$. \Cref{fig:search_superdiffusion}\textbf{b,d)} shows these fitted exponents $\lambda(q)$ against $q$. The linear fit is good, with $\lambda(q) = (0.897 \pm 0.005)\,q$ and $R^2 = 0.999$, across the whole considered range $0.1 \le q \le 3.0$. This gives a Lévy index $\alpha = 1.115 \pm 0.007$, which is well below the Brownian value of $2$ and so indicates strong superdiffusion. Because a linear regression well describes the spectrum, there is no crossover in scaling between the low-$q$ moments that probe fixational motion and the high-$q$ moments that probe saccades. A single exponent describes both, validating our unifying model across both scales of study in gaze dynamics, derived naturally from FNS dynamics.
\begin{figure}[htbp!]
	\centering
	\begin{subfigure}[t]{0.49\textwidth}
		\centering
		\panel{a)}{figures/superdiffusion/data/rescaled_moment_spectrum_data.pdf}
		\label{subfig:fad_data}
	\end{subfigure}
	\hfill
	\begin{subfigure}[t]{0.49\textwidth}
		\centering
		\panel{b)}{figures/superdiffusion/data/lambda_vs_q_fit_data.pdf}
		\label{subfig:scaling_order_data}
	\end{subfigure}
	\vfill
	\begin{subfigure}[t]{0.49\textwidth}
		\centering
		\panel{c)}{figures/superdiffusion/model/rescaled_moment_spectrum.pdf}
		\label{subfig:fad_model}
	\end{subfigure}
	\hfill
	\begin{subfigure}[t]{0.49\textwidth}
		\centering
		\panel{d)}{figures/superdiffusion/model/lambda_vs_q_fit.pdf}
		\label{subfig:scaling_order_model}
	\end{subfigure}
    \caption{\textbf{The model reproduces the monoscaling superdiffusion of human visual search through a linear spectrum of moments.} Rescaled moments $M_q(\Delta t)$ (\cref{eq:Mq}) and fitted scaling exponents $\lambda(q)$ for 30 values of $q$ spaced linearly from $0.1$ to $3.0$. \textbf{Top row:} binocular-averaged gaze data from the 22 subjects of \textcite{naher_primate_saccades} (1846 trajectories, mean duration \qty{5}{\s}). \textbf{Bottom row:} 500 trajectories simulated from the model of \cref{chap:a_unifying_model} (mean duration \qty{15}{\s}). \textbf{(a,c)}~Moments on log--log axes. The shaded region denotes the fitting regime, and the dashed line shows the $q=2$ fit, which recovers the square-root scaling of the traditional MSD. Ringing oscillations in the low--$q$ moments result from the highly regular harmonic oscillations that model gamma-scale OMT. \textbf{(b,d)}~The spectra are linear, $\lambda(q) = q/\alpha$, in both cases: $\alpha = 1.115 \pm 0.007$ for the data ($R^2 = 0.999$) and $\alpha = 1.091 \pm 0.005$ for the model, showing robust superdiffusion across scales, compared with $\alpha = 2$ recovering Brownian motion.}
    \label{fig:search_superdiffusion}
\end{figure}
\subsubsection*{Implications of monoscaling superdiffusion}
The most important implication of this particular superdiffusion result is that the strong monoscaling of visual search trajectories with Lévy index $\alpha = 1.115$ is robustly multiscale. The robust linear scaling of $\lambda(q)=\frac{q}{\alpha}=\frac{q}{1.115}$ across a spectrum of $q$-values $0 < q < 3$ is important, as it suggests an identical process driving both fixational eye movements --- including tremor, drifts, and microsaccades --- with small relative values of $\| \Delta \textbf{x}\|$ emphasised with low values of $q$, as well as large nonlocal saccades with large relative values of $\| \Delta \textbf{x} \|$, emphasised through large values of $q$. Importantly, this validates our construction of a continuous-time unified model of eye movements, unified across all spatiotemporal scales of visual search.
